\einheitenkopf{Einheit 5 von 5 · Diagramme beurteilen und Boxplot}
\erg{36}{a) ja \quad b) nein (880) \quad c) ja \quad d) nein (840) \quad e) ja \quad f) nein -- von 2021 auf 2022 gehen die Ausleihen von 910 auf 880 zurück \quad g) nein: $2 \cdot 840 = 1680$, 2024 waren es nur 1200 \quad h) ja: $840 : 4 = 210$, $840 + 210 = 1050$ \quad i) nein: $1 : 8 = 0{,}125 = 12{,}5\,\%$ \quad j) Beide Teile sind falsch: der Anstieg war nicht durchgehend (Rückgang 2022), und verdoppelt hätte $1680$ bedeutet, erreicht wurden 1200.}
\erg{37}{a) 400 \quad b) $(450-410) : 410 = 0{,}0975\ldots \approx 9{,}8\,\%$ \quad c) Die Achse beginnt bei 400; jede Säule zeigt nur den Teil über 400. Die Höhen verhalten sich wie $10 : 20 : 30 : 50$, die Mitgliederzahlen aber nur wie $410 : 420 : 430 : 450$. \quad d) Mit einer Achse ab 0 sind alle vier Säulen fast gleich hoch, der Zuwachs ist kaum noch zu sehen. \quad e) Die Fortschreibung ist unzulässig: Aus vier Jahren mit $+40$ Mitgliedern folgt nicht, dass es so weitergeht; nach diesem Muster wären es 2030 etwa 530, nicht über 1000. \quad f) Die Überschrift trifft nicht zu. Tatsächlich sind es 40 Mitglieder mehr, also rund $9{,}8\,\%$ in drei Jahren. Der Eindruck entsteht allein durch die bei 400 abgeschnittene Achse.}
\erg{38}{a) 3 \quad b) 19 \quad c) 11 \quad d) 7 \quad e) 15 \quad f) $19-3=16$ \quad g) $50\,\%$ \quad h) $25\,\%$}
\erg{39}{a) Tarek hat aus der doppelten Säulenhöhe auf den doppelten Wert geschlossen. Weil die Achse erst bei 200 beginnt, zeigt die Säule nur den Teil über 200; doppelte Höhe heißt doppelter Anteil über 200, nicht doppelter Wert. \quad b) $(240-220) : 220 = 0{,}0909\ldots \approx 9{,}1\,\%$}
\erg{40}{a) Nur bei einer Achse ab 0 ist die Säulenhöhe zum Wert proportional. Sonst zeigt die Höhe nur den Teil über dem Achsenanfang, und Höhenverhältnisse führen in die Irre. \quad b) Die Behauptung ist nicht begründet. Das Diagramm enthält nur Werte bis 2024; was danach geschieht, steht dort nicht. Ein Rückgang muss nicht gleichmäßig weitergehen und endet nicht zwangsläufig bei null.}
